% ============================================================
%  FORMELSAMMLUNG-TEMPLATE
%  Dreispaltiges A4-Layout für Klausur-Formelsammlungen
%
%  DATEISTRUKTUR:
%    TEMPLATE.tex           ← diese Datei (Preamble + \input-Liste)
%    parts/01_kapitel.tex   ← Kapitel-Dateien (je Vorlesungsskript eine)
%    parts/02_kapitel.tex
%    ...
%
%  NEUE KAPITEL HINZUFÜGEN:
%    1. Neue Datei in parts/ anlegen
%    2. \input{parts/XX_name} am Ende dieser Datei eintragen
%    3. Eigene Farbe in Farbschema definieren (cMeinKapitel)
% ============================================================
\documentclass[a4paper,10pt]{article}
\usepackage[utf8]{inputenc}
\usepackage[T1]{fontenc}
\usepackage[ngerman]{babel}
\usepackage{amsmath,amssymb,mathtools}
\usepackage{xcolor}
\usepackage{geometry}
\usepackage{array}
\usepackage{multirow}
\usepackage{microtype}
\geometry{top=10mm,bottom=10mm,left=10mm,right=10mm}
\setlength{\parskip}{0pt}
\setlength{\parindent}{0pt}
\renewcommand{\baselinestretch}{0.9}
\setlength{\emergencystretch}{2em}
\setlength{\fboxsep}{0.18em}
\setlength{\fboxrule}{0.3pt}
\hfuzz=3pt   % dfrac in engen Spalten: <3pt Overfull unterdrücken

% ============================================================
%  FARBSCHEMA — je Kapitel/Skript eine Farbe
%  Syntax: \colorlet{cMeinName}{Farbe}
%  Verfügbare Basisfarben: blue, teal, violet, orange, red,
%                          green!60!black, cyan!70!black, ...
% ============================================================
\colorlet{cKap1}{blue}            % Teil 1–4: Signale, Systeme, Stabilität, Nichtlinearität
\colorlet{cKap2}{teal}            % Variante 2
\colorlet{cKap3}{violet}          % Variante 3
\colorlet{cKap4}{orange!90!red}   % Variante 4

% ============================================================
%  MAKROS — ÜBERSICHT
%
%  \bereich[Farbe]{Titel}      Farbiger Abschnittsbalken
%  \bereichende                Trennlinie zwischen Abschnitten
%  formelreihe (Umgebung)      3-spaltiger Tabellenblock
%  \formelblock{Name}{Formel}  Formel mit Bezeichnung (grau/kursiv)
%  \beispielblock{Name}{Inh.} Hervorgehobener Beispielblock (orange)
%  \infoblock{Text}            Kursiver Erklärungstext in Spalte
% ============================================================

% ---- Abschnittsbalken: \bereich[cKap1]{Titel} --------------
% Spart ~2 Zeilen gegenüber \section*
% Standardfarbe: cKap1 (kann weggelassen werden)
\newcommand{\bereich}[2][cKap1]{%
  \vspace{0.5em}%
  \noindent\colorbox{#1!10}{%
    \makebox[\dimexpr\linewidth-2\fboxsep\relax][l]{%
      \small\bfseries\color{#1!80!black}\strut\enspace #2}}%
  \vspace{0.1em}\par\noindent%
}

% ---- Trennlinie --------------------------------------------
\newcommand{\bereichende}{%
  \vspace{0.2em}%
  {\color{gray!50}\hrule height 0.3pt}%
  \vspace{0.2em}%
}

% ---- 3-Spalten-Layout --------------------------------------
% Jede Zeile endet mit \\ oder \\[Abstand]
% Spalten werden mit & getrennt
% Leerzeile zwischen Zeilengruppen: & & \\[0.3em]
\newenvironment{formelreihe}{%
  \noindent\footnotesize%
  \setlength{\tabcolsep}{0pt}%
  \renewcommand{\arraystretch}{0.85}%
  \begin{tabular}{@{}p{0.32\linewidth}@{\hspace{0.02\linewidth}}%
                     p{0.32\linewidth}@{\hspace{0.02\linewidth}}%
                     p{0.32\linewidth}@{}}%
}{%
  \end{tabular}\par%
}

% ---- Formel-Block ------------------------------------------
% #1 = Bezeichnung (erscheint in Kapitälchen grau über der Formel)
% #2 = LaTeX-Mathe (wird automatisch in \displaystyle gesetzt)
% Für mehrzeilige Formeln: aligned-Umgebung in #2 verwenden
\newcommand{\formelblock}[2]{%
  \parbox[t]{\linewidth}{%
    {\scriptsize\scshape\color{gray!65!black} #1}\\[0.03em]%
    $\displaystyle #2$%
  }%
}

% ---- Beispiel-Block ----------------------------------------
% Orangefarbener Rahmen — immer in Spalte 3 (rechte Spalte)
% #1 = kurze Kontextbeschreibung
% #2 = Inhalt (Mathe oder Text, wird in \displaystyle gesetzt)
% Für Zeilenumbruch in #2: aligned-Umgebung verwenden
\newcommand{\beispielblock}[2]{%
  \fcolorbox{orange!50}{yellow!15}{%
    \parbox[t]{\dimexpr\linewidth-2\fboxsep-2\fboxrule\relax}{%
      {\scriptsize\scshape\color{orange!80!black} Bsp: #1}\\[0.03em]%
      $\displaystyle #2$%
    }}%
}

% ---- Info-Block --------------------------------------------
% Kursiver Fließtext — für Definitionen/Kontext in Spalte 1
\newcommand{\infoblock}[1]{%
  \parbox[t]{\linewidth}{\scriptsize\itshape #1}%
}

% ============================================================
\begin{document}
\begin{center}
  {\large\bfseries Formelsammlung — SOS/A (Nichtlineare und Stochastische Systeme)}
\end{center}
\vspace{0.3em}

% ---- Kapitel einbinden -------------------------------------
% Auskommentierte Kapitel werden ignoriert (nützlich beim Aufbau)
% ============================================================
%  SOS/A FORMELSAMMLUNG — Teil 1–4
%  Signale, Systeme, Stabilitätsanalyse, Nichtlinearität
% ============================================================

\bereich[cKap1]{Teil 1: Signale, Systeme und deren Eigenschaften}

\begin{formelreihe}
  \infoblock{Ein Signal ist eine Funktion mindestens einer unabhängigen Variable, die Information repräsentiert.} &
  \infoblock{Ein System ist eine Abbildung: \(u(t) \to y(t)\) oder \(\mathbf{u}(t) \to \mathbf{y}(t)\)} &
  \\[0.3em]

  \formelblock{Lineare Operation}{\begin{aligned}
    y(t) &= c \cdot u(t) \\
    y(t) &= \frac{\mathrm{d}u(t)}{\mathrm{d}t} \\
    y(t) &= \int_{0}^{t} u(\tau) \, \mathrm{d}\tau
  \end{aligned}} &
  \formelblock{Nichtlineare Operation}{\begin{aligned}
    y(t) &= u^2(t) \\
    y(t) &= u(t) + k \quad (k\neq 0) \\
    y(t) &= \sin(u(t))
  \end{aligned}} &
  \beispielblock{Linear}{
    \(y(t)=5u(t)\)\\[0.1em]
    konstanter Faktor} \\[0.3em]

  \formelblock{Statisches System}{Gedächtnislos:\\
    \(y(t) = f(u(t))\)} &
  \formelblock{Dynamisches System}{Mit Gedächtnis:\\
    \(y(t) = f(u(t), u(t-1), \ldots, x)\)} &
  \beispielblock{Statisch}{Verstärker:\\
    \(y=Ku\)} \\[0.3em]

  \infoblock{Zeitinvariant:\\System-verhalten ändert sich\\nicht mit Zeit.\\
    \(y(t-T) \Rightarrow y_{\text{delay}}(t) = y(t-T)\)} &
  \infoblock{Zeitvariant:\\Systemparameter\\ändern sich mit Zeit.\\
    z.\,B. Funkkanal, Satellit} &
  \\
\end{formelreihe}
\bereichende

% ============================================================

\bereich[cKap2]{Teil 2: Stabilitätsuntersuchung von LZI-Systemen — Routh-Schema}

\begin{formelreihe}
  \formelblock{BIBO-Stabilität}{Ein LZI-System mit Übertragungsfunktion \(G(s)\) ist BIBO-stabil \(\Leftrightarrow\) alle Pole haben negativen Realteil.} &
  \infoblock{Alle Pole der Übertragungs-funktion \(G(s) = \frac{N(s)}{Z(s)}\) liegen in der linken s-Ebene.} &
  \beispielblock{Pole}{
    \(G(s)=\frac{1}{s+1}\)\\
    Pol: \(s=-1\)\\
    \(\Rightarrow\) stabil} \\[0.3em]

  \formelblock{Routh-Schema Regel 1}{Alle Koeffizienten von \(N(s)\) müssen gleiches Vorzeichen haben.\\
    \(N(s) = a_ns^n + a_{n-1}s^{n-1} + \cdots + a_0\)\\
    Wenn nicht: \textbf{nicht stabil}} &
  \infoblock{Notwendig aber nicht hinreichend! Ausnahme: \(n=1, 2\) — dann auch hinreichend.} &
  \beispielblock{Prüfung}{
    \(N(s)=s^2+3s+2\)\\
    Alle \(>\,0\) \(\Rightarrow\)\\könnte stabil sein} \\[0.3em]

  \formelblock{Routh-Schema Regel 2}{Zähle Vorzeichenwechsel in linker Spalte des Schemas.\\
    Anzahl Wechsel = Anzahl Pole rechts} &
  \formelblock{Routh-Schema Regel 3}{Null in linker Spalte\\
    \(\Rightarrow\) mindestens ein Pol auf imaginärer Achse} &
  \beispielblock{Beispiel}{
    \(N(s)=s+1\) (n=1)\\
    Koeff: \(1, 1\)\\
    Alle \(>\,0\) \(\Rightarrow\) stabil} \\[0.3em]

  \infoblock{Für \(n=1\): \(N(s)=as+b\)\\
    \(\Rightarrow\) stabil wenn \(a,b\) gleiches Vorzeichen} &
  \infoblock{Für \(n=2\): \(N(s)=as^2+bs+c\)\\
    \(\Rightarrow\) stabil wenn alle gleich\\positiv oder alle gleich negativ} &
  \beispielblock{n=2}{
    \(N=s^2+4s+1\)\\
    Alle \(>\,0\) \(\Rightarrow\) stabil} \\
\end{formelreihe}
\bereichende

% ============================================================

\bereich[cKap3]{Teil 3: Zustandsraumdarstellung und Übertragungsfunktion}

\begin{formelreihe}
  \formelblock{Zustandsraum (allgemein)}{\begin{aligned}
    \dot{\mathbf{x}}(t) &= \mathbf{A}\,\mathbf{x}(t) + \mathbf{B}\,u(t) \\
    y(t) &= \mathbf{C}\,\mathbf{x}(t) + D\,u(t)
  \end{aligned}} &
  \infoblock{\(\mathbf{A}\): System-matrix, \(\mathbf{B}\): Eingangs-matrix, \(\mathbf{C}\): Ausgangs-matrix, \(D\): Durchgriff} &
  \beispielblock{Dyn. System}{RC-Glied:\\
    \(\dot{x}=-\frac{1}{RC}x + \frac{1}{RC}u\)\\
    \(y = x\)} \\[0.3em]

  \formelblock{Übertragungsfunktion aus ZR}{\begin{aligned}
    G(s) &= \mathbf{C}(s\mathbf{I}-\mathbf{A})^{-1}\mathbf{B} + D
  \end{aligned}} &
  \formelblock{Charakterist. Polynom}{\begin{aligned}
    N(s) &= \det(s\mathbf{I}-\mathbf{A}) \\
    &= \text{Nennerpolynom}
  \end{aligned}} &
  \beispielblock{Eigenwerte}{Pole von \(G(s)\)\\
    sind Eigenwerte\\von \(\mathbf{A}\)} \\[0.3em]

  \infoblock{Faltung (Zeit):\\
    \(y(t) = \int_0^t h(\tau)\,u(t-\tau)\mathrm{d}\tau\)} &
  \formelblock{Multiplikation (Freq.)}{
    \(Y(s) = G(s) \, U(s)\)\\
    \(G(s) = \frac{Y(s)}{U(s)}\)} &
  \beispielblock{1. Ordnung}{\(G(s)=\frac{K}{Ts+1}\)\\
    \(\Rightarrow\)\\
    \(\dot{x}=-\frac{1}{T}x+\frac{K}{T}u\)} \\
\end{formelreihe}
\bereichende

% ============================================================

\bereich[cKap4]{Teil 4: Nichtlineare Systeme — Ruhelagen \& Stabilitätsanalyse}

\begin{formelreihe}
  \formelblock{Nichtlin. System (allg.)}{\begin{aligned}
    \dot{\mathbf{x}}(t) &= \mathbf{f}(\mathbf{x}(t), u(t), t) \\
    y(t) &= h(\mathbf{x}(t), u(t), t)
  \end{aligned}} &
  \infoblock{Nichtlineare Systeme können einzelne isolierte Ruhelagen haben, chaotisches Verhalten zeigen, Hysterese-Effekte aufweisen.} &
  \beispielblock{Pendel}{\begin{aligned}
    \dot{x}_1 &= x_2\\
    \dot{x}_2 &= -\frac{g}{\ell}\sin(x_1)
  \end{aligned}} \\[0.3em]

  \formelblock{Ruhelage \(\mathbf{x}_R\)}{Autonomes System (\(u(t)=0\)):\\
    \(\mathbf{f}(\mathbf{x}_R, \mathbf{0}) = \mathbf{0}\)} &
  \formelblock{Betriebspunkt \(\mathbf{x}_B\)}{Mit Eingangssignal:\\
    \(\mathbf{f}(\mathbf{x}_B, u_B) = \mathbf{0}\)\\
    bei \(u(t)=u_B \, \forall t\)} &
  \beispielblock{Pendel-Ruhelagen}{Hängend: \(x_1=0\)\\
    Stehend: \(x_1=\pi\)\\
    \(\sin(x_1)=0\)} \\[0.3em]

  \formelblock{Vektornormen}{\(\|\mathbf{x}\|_1 = \sum_i |x_i|\)\\
    \(\|\mathbf{x}\|_2 = \sqrt{\sum_i x_i^2}\)\\
    \(\|\mathbf{x}\|_\infty = \max_i |x_i|\)} &
  \infoblock{Norm-Eigenschaften: \(\|\mathbf{x}\| = 0 \Leftrightarrow \mathbf{x}=\mathbf{0}\), \(\|c\,\mathbf{x}\| = |c|\,\|\mathbf{x}\|\), Dreiecksungleichung.} &
  \beispielblock{2-Norm}{
    \(\mathbf{x}=\begin{bmatrix}3\\4\end{bmatrix}\)\\
    \(\|\mathbf{x}\|_2=5\)} \\[0.3em]

  \formelblock{Stabil nach Lyapunov}{\(\forall \varepsilon > 0 \, \exists \delta > 0\):\\
    wenn \(\|\mathbf{x}(0)-\mathbf{x}_R\| < \delta\)\\
    dann \(\|\mathbf{x}(t)-\mathbf{x}_R\| < \varepsilon \, \forall t \geq 0\)} &
  \formelblock{Asymptotisch stabil}{Stabil PLUS:\\
    \(\lim_{t\to\infty} \|\mathbf{x}(t)-\mathbf{x}_R\| = 0\)\\
    für \(\mathbf{x}(0)\) nahe \(\mathbf{x}_R\)} &
  \beispielblock{Hängendes Pendel}{Mit Reibung:\\
    Asymptotisch stabil} \\[0.3em]

  \formelblock{Global asympt. stabil}{Asymptotisch stabil für beliebige Anfangsbedingungen \(\mathbf{x}(0) \in \mathbb{R}^n\).\\
    Unbegrenztes Einzugsgebiet.} &
  \infoblock{Übergangsdynamik unabhängig von Anfangsbedingung. System konvergiert immer.} &
  \beispielblock{PT1-Glied}{\(G(s)=\frac{K}{Ts+1}\)\\
    Global asymptotisch stabil} \\
\end{formelreihe}
\bereichende

% ============================================================

\bereich[cKap4]{Linearisierung nichtlinearer Systeme im Betriebspunkt}

\begin{formelreihe}
  \formelblock{Taylor-Linearisierung}{\(\dot{x}_i = f_i(\mathbf{x}, u)\)\\
    Um \((\mathbf{x}_B, u_B)\):\\
    \(\Delta\dot{x}_i \approx \sum_j \frac{\partial f_i}{\partial x_j}\Delta x_j + \frac{\partial f_i}{\partial u}\Delta u\)} &
  \formelblock{In Matrixform}{\(\Delta\dot{\mathbf{x}} = \mathbf{A}\,\Delta\mathbf{x} + \mathbf{B}\,\Delta u\)\\
    mit \(\mathbf{A} = \frac{\partial \mathbf{f}}{\partial \mathbf{x}}\),\\
    \(\mathbf{B} = \frac{\partial \mathbf{f}}{\partial u}\) (bei \((\mathbf{x}_B, u_B)\))} &
  \beispielblock{Pendel}{
    \(f = \sin(x_1)\)\\
    \(\frac{\partial f}{\partial x_1}= \cos(x_1)\)\\
    Bei \(x_1=0\): Steigung 1} \\[0.3em]

  \formelblock{Delta-Koordinaten}{Transformation:\\
    \(\mathbf{x} = \Delta\mathbf{x} + \mathbf{x}_B\)\\
    \(u = \Delta u + u_B\)\\
    Betriebspunkt wird zur\\Ruhelage des \(\Delta\)-Systems} &
  \infoblock{Vorteil: Linearisiertes System hat Ruhelage im Ursprung \(\Delta\mathbf{x}=0\). Lyapunov und BIBO-Kriterien anwendbar.} &
  \beispielblock{Beispiel}{\(\dot{x}=u-x^3+2\)\\
    \(u_B=2 \Rightarrow x_B=0\)\\
    Linearisiert:\\
    \(\Delta\dot{x} = \Delta u\)} \\
\end{formelreihe}
\bereichende

%\input{parts/02_kapitel}
%\input{parts/03_kapitel}
%\input{parts/04_kapitel}

\end{document}
